\documentclass[11pt,a4paper]{article}

\usepackage[margin=2.5cm]{geometry}
\usepackage{hyperref}
\usepackage{booktabs}
\usepackage{amsmath}
\usepackage{parskip}
\usepackage{microtype}
\usepackage{enumitem}
\usepackage{xcolor}
\usepackage{titlesec}

\hypersetup{colorlinks=true, linkcolor=blue, urlcolor=blue}

\definecolor{commitgray}{rgb}{0.35,0.35,0.35}
\newcommand{\commit}[1]{\texttt{\textcolor{commitgray}{#1}}}

\titleformat{\section}{\large\bfseries}{}{0em}{}[\titlerule]
\titleformat{\subsection}{\normalsize\bfseries}{}{0em}{}

\title{\textbf{LoRA4VR --- Training Log}\\
\normalsize A living record of experiments, findings, and open questions}
\author{}
\date{}

\begin{document}
\maketitle

\noindent\textbf{What this is.}
A dated log of experiments run on the LoRA4VR self-play loop. Each entry
records what was tried, what happened, and how to reproduce it. Entries are
added as experiments complete. Reproduction references a git commit; the code
at that commit is the exact code that produced the result.

\noindent\textbf{System.} \texttt{Qwen2.5-Coder-3B-Instruct} fine-tuned with
GRPO/DAPO on a Google Cloud L4 GPU (24\,GB). The loop generates
\texttt{policy(state) -> action} Python code, scores it by running the policy
in a sandboxed environment, and trains LoRA adapters on the reward signal.

\bigskip
\hrule
\bigskip

% ============================================================
\section*{2026-04-13 \quad Environment design and first training runs}
% ============================================================

\commit{d96f7cb} $\cdot$ \textit{puck-on-tilted-plane env, runs 1--3}

\subsection*{Motivation}

CartPole was the original seed task but it is trivially solved by a two-line
bang-bang policy, giving a flat learning curve from iteration 0. Any standard
RL benchmark (LunarLander, Pendulum, MountainCar) is heavily represented in
the base model's pretraining data, so improvement from LoRA would be
confounded with prior knowledge. We need an environment the model has to
\emph{read} rather than recall.

\subsection*{The puck environment}

A puck moves in a bounded $[-6, 6]^2$ arena under a constant hidden tilt
$(t_x, t_y)$ and linear drag $\mu$. The policy receives
$\{p_x, p_y, v_x, v_y, g_x, g_y\}$ and must drive the puck to the goal
within 120 steps. Hidden parameters --- tilt, drag, thrust limit $a_{\max}$,
initial velocity --- are sampled per scenario from a seed and are not in the
state dict. The tight thrust envelope ($a_{\max} = \|tilt\|_2 \times 1.05 +
\epsilon$) ensures that a naive ``thrust straight at the goal'' policy
fails because it cannot cancel tilt drift.

\textbf{Reward:} $-\Delta t \cdot \text{dist}(p, g)$ per step; $+50$ bonus
on reaching the goal. Errored episodes score $-200$ per scenario (floor).

\textbf{Baseline gap} (200 seeds, hand-written policies):
\begin{itemize}[nosep]
  \item Naive aim: $\approx 26\%$ goal rate
  \item PD + tilt feedforward: $\approx 73\%$ goal rate
\end{itemize}

This gap is real headroom for LoRA to exploit.

\subsection*{Bugs encountered and fixed}

Three bugs had to be resolved before any learning was possible.

\textbf{Run 1 --- OOM on L4.}
Default config ($G=16$, batch 16, max completion 900 tokens) ran out of
memory at the backward pass. Fix (PR \#6): halved $G$ and batch to 8,
trimmed completion length to 600.

\textbf{Run 2 --- Reward sign flip.}
Training ran but degraded (fewer running samples each iteration). Errored
samples scored $0.0$ while samples that ran and drifted integrated large
negative rewards ($\approx -280$). GRPO's normalized advantage placed the
$0.0$ errors \emph{above} the mean, so the trainer rewarded broken code.
Fix (PR \#7): errored scenarios score $-200$, making broken code strictly
the worst outcome.

\textbf{Run 3 --- Dead advantages.}
With LR $= 10^{-5}$ and only 1 training prompt (4 gradient steps per outer
iteration), $\approx 75\%$ of batches were uniformly errored
(advantage $= 0$, no gradient). On the remaining steps, weight deltas were
$\sim 10^{-6}$ --- indistinguishable from zero. The reward was identical to
7 decimal places at iteration 0 and iteration 9. Fix (PR \#10): LR bumped
$10\times$ to $10^{-4}$; epochs raised to 20 (20 gradient steps per outer
iteration on 1 prompt).

\subsection*{Reproduce}

Commits: OOM fix \commit{ef228fe}, reward floor \commit{ef7e0c9},
LR/epoch fix \commit{e9353d7}. Env source: \texttt{scripts/seed\_tasks/puck\_tilted\_plane\_env.py}.

\bigskip
\hrule
\bigskip

% ============================================================
\section*{2026-04-18 \quad Run 5 probe: base model vs.\ 5 iterations of GRPO}
% ============================================================

\commit{d5bc81c} $\cdot$ \textit{puck\_grpo\_run\_5, iters 0 and 5, $n=20$ samples each}

\subsection*{Setup}

Run 5 used the config at \commit{e9353d7} (LR $= 10^{-4}$, $G=8$, 20 epochs
per outer iteration, DAPO loss, $\beta = 0$). It completed 5 outer iterations
before the VM was stopped. Adapter checkpoints were recovered to
\texttt{models/lora/puck\_grpo\_run\_5/}.

To inspect what the model actually writes, we sample 20 policies from the
base model (iter 0, no adapters) and from the trained model (iter 5), score
each against 8 scenarios, and compare. The tool is
\texttt{scripts/probe\_generation.py}; it saves per-sample policy files,
a summary CSV, and per-scenario reward breakdown to
\texttt{data/probe/\{run\_name\}\_iter\{N\}\_t\{temp\}/}.

\subsection*{Results}

\begin{table}[h]
\centering
\begin{tabular}{lrrrr}
\toprule
 & \textbf{Ran OK} & \textbf{Crashed} & \textbf{Mean reward (runners)} & \textbf{Best reward} \\
\midrule
Iter 0 (base) & 13 / 20 & 7 / 20 & $-218.9$ & $-92.4$ \\
Iter 5        & 20 / 20 & 0 / 20 & $-140.4$ & $-140.4$ \\
\bottomrule
\end{tabular}
\end{table}

\subsubsection*{Iter 0}

The base model generates varied PD/PID controllers but crashes 35\% of the
time by attempting to read hidden env internals (\texttt{a\_max}, \texttt{dt},
\texttt{steps}) from the state dict. The 13 successful policies span a reward
range of $[-92, -373]$ with genuine variation in gain choices. The best
policy uses asymmetric gains ($k_{px}=2.0, k_{dx}=1.0$; $k_{py}=1.5,
k_{dy}=1.0$) and scores $-92.4$ over 8 scenarios.

\subsubsection*{Iter 5}

Zero crashes. The error-reading habit is gone. However, all 20 samples at
temperature $0.8$ produce \emph{exactly} $-140.3516$ (identical to 4 decimal
places). Every policy is byte-for-byte the same: $k_p = 2.0$, $k_d = 1.5$.
This is the \textbf{illustrative example embedded in the prompt template}
(\texttt{configs/prompt\_template.txt}). The model has memorized and
reproduced it verbatim.

\subsection*{Interpretation}

Training achieved two things and failed at a third.

\textit{Achieved.} Error rate dropped from 35\% to 0\% --- the model learned
not to read hidden env attributes. Mean reward over runners improved by 35\%.

\textit{Failed.} The model collapsed to a single policy. With $\beta = 0$
(no KL penalty), one training prompt, and 20 gradient steps per outer
iteration, GRPO found the prompt example as a high-reward fixed point and
reinforced it into a delta distribution. By iter 5, temperature $0.8$
sampling produces zero diversity. The best reachable policy is now $-140.4$,
worse than the base model's best ($-92.4$).

This is reward hacking: the model optimized the training distribution
(memorize the example) rather than the environment.

\subsection*{Hypotheses to test}

\begin{enumerate}[nosep]
  \item \textbf{Remove or paraphrase the prompt example.} It is the attractor
        GRPO is collapsing to. A deliberately suboptimal or absent example
        removes the fixed point.
  \item \textbf{More training prompts.} With 1 prompt, every outer iteration
        trains on the same input. Adding 4--8 tasks with varied dynamics forces
        generalization.
  \item \textbf{Raise $\beta > 0$.} A non-zero KL coefficient prevents the
        policy from drifting too far from the base distribution and preserves
        diversity.
  \item \textbf{Reduce epochs.} 20 gradient steps on 1 prompt per outer
        iteration is aggressive. 5--10 may preserve more exploration.
  \item \textbf{Probe at every intermediate checkpoint.} Does collapse happen
        at iter 1 or gradually? One probe per checkpoint would pinpoint when
        diversity dies.
\end{enumerate}

\subsection*{Open questions}

\begin{itemize}[nosep]
  \item What does the W\&B reward curve look like across the 5 iterations?
        (The run at \texttt{puck\_grpo\_run\_5-orchestrator} has data but some
        post-resume steps may be missing due to a step-monotonicity issue.)
  \item What is the PD+feedforward ceiling expressed in probe-script reward
        units? This would establish what score a successful LoRA policy needs
        to beat.
\end{itemize}

\subsection*{Reproduce}

At commit \commit{d5bc81c}, with adapter checkpoints in
\texttt{models/lora/puck\_grpo\_run\_5/}:

\smallskip
\texttt{scripts/probe\_generation.py} \texttt{--iter 0 --run\_name puck\_grpo\_run\_5 --n 20}

\texttt{scripts/probe\_generation.py} \texttt{--iter 5 --run\_name puck\_grpo\_run\_5 --n 20}

\smallskip
Results land in \texttt{data/probe/puck\_grpo\_run\_5\_iter\{000,005\}\_t0p80/}.

\end{document}
